%%%%%%%%%%%%%%%%
%%% num 6 %%%
%%%%%%%%%%%%%%%%

\Chapter{6}

\Verse{1}
{
	\hP{וַיְדַבֵּר יְהֹוָה אֶל מֹשֶׁה לֵּאמֹר׃}
}
{
	\eP{6 And YHWH spoke to Moses, saying, [image]}
}
{
}

\Verse{2}
{
	\hP{דַּבֵּר אֶל בְּנֵי יִשְׂרָאֵל וְאָמַרְתָּ אֲלֵהֶם אִישׁ אוֹ אִשָּׁה כִּי יַפְלִא לִנְדֹּר נֶדֶר נָזִיר לְהַזִּיר לַיהֹוָה׃}
}
{
	\eP{
		“Speak to the children of Israel, and you shall say to\\
		them: When a man or a woman will expressly make a Nazirite\\
		vow to make a separation for YHWH, {[image]}\\
	}
}
{
%	\footnote{
%		a man or a woman. Although the rest of the law is stated in the
%		masculine, the entire matter of choosing to be a Nazirite applies to
%		both men and women. Footnote 6:2. make a separation for YHWH. A Nazirite
%		vows to give up certain things and live in a state of holiness. There is
%		the general idea in the Torah that all the people are holy (Lev 19:1),
%		but the Nazirite vow refers to a singular state that exceeds that. The
%		Nazirite gives up three things: consuming alcohol, coming near to dead
%		persons, and cutting his or her hair. Why these three things? These
%		three things are related to priesthood—and particularly to the high
%		priest. Israel’s priests are forbidden to consume alcohol at the Tent of
%		Meeting, and this is the first command in the Torah that God gives
%		directly to Aaron, the high priest, not to Moses (Lev 10:9). Priests are
%		forbidden to come to dead persons except their closest relatives (Lev
%		21:1–3), and the high priest is commanded that “he shall not come to any
%		dead persons; he shall not become impure for his father and for his
%		mother” (21:11)—which is the wording here for the Nazirite as well: “he
%		shall not come to a dead person. For his father and for his mother, for
%		his brother and for his sister: he shall not become impure for them”
%		(Num 6:6–7). But in the matter of the hair, the high priest is commanded
%		to do the opposite of the Nazirite: The Nazirite shall “grow the hair of
%		his head loose” (Num 6:5), but the high priest “shall not let loose the
%		hair of his head” (Lev 21:10). What does it mean that a Nazirite is
%		someone who chooses to live in certain respects like a high priest but
%		that a high priest is commanded not to have the most obvious
%		characteristic of the Nazirite? The issue appears to me to relate to the
%		fact that the clergy in Israel is not open to most Israelites to choose,
%		but only to members of the tribe of Levi by heredity, and then only to
%		males; and the priesthood comes to be limited only to a certain family
%		within Levi, the family of Aaron; and the high priesthood is limited to
%		a son of the preceding high priest. Prophets can be from any tribe and
%		can be male or female; one becomes a prophet by being called by God.
%		Postbiblically, rabbis replace priests as the primary clergy of the
%		Jews, and membership in the rabbinate is open to anyone who chooses it
%		and proves capable of qualifying and carrying it out. Originally it was
%		limited to males, but in the present generation the Reform,
%		Conservative, and Reconstructionist movements have included both females
%		and males. So what option does a person in the biblical world have who
%		is drawn to the holy life of a priest but who is not a male Levite from
%		the family of Aaron? Such a person can choose to be a Nazirite. And note
%		the case of Samuel: he becomes a priest, seemingly the high priest, even
%		though he is not a Levite—and Samuel is a Nazirite! He replaces the
%		Levite high priest Eli at the Tabernacle at Shiloh, and he performs
%		sacrifices, which no nonpriest, even the king, may do (1 Samuel 1–13).
%		Historically, the law forbidding the high priest to let his hair grow
%		loose may have been composed precisely to establish that this could
%		never happen again. No one outside the Aaronid Levite clan could just
%		choose to enter the clergy and rise in the priesthood. As the law stands
%		in the Torah now, preceding the story of Samuel, it means that the case
%		of Samuel must be regarded as unique, with Samuel’s rise in the
%		priesthood being the result of a special selection by God.
%	}
}

\Verse{3}
{
	\hP{מִיַּיִן וְשֵׁכָר יַזִּיר חֹמֶץ יַיִן וְחֹמֶץ שֵׁכָר לֹא יִשְׁתֶּה וְכל מִשְׁרַת עֲנָבִים לֹא יִשְׁתֶּה וַעֲנָבִים לַחִים וִיבֵשִׁים לֹא יֹאכֵל׃}
}
{
	\eP{
		he shall separate from wine and beer; he shall not drink\\
		any vinegar of wine or vinegar of beer, and he shall not\\
		drink any juice of grapes, and he shall not eat fresh or\\
		dried grapes. {[image]}\\
	}
}
{
}

\Verse{4}
{
	\hP{כֹּל יְמֵי נִזְרוֹ מִכֹּל אֲשֶׁר יֵעָשֶׂה מִגֶּפֶן הַיַּיִן מֵחַרְצַנִּים וְעַד זָג לֹא יֹאכֵל׃}
}
{
	\eP{
		All the days of his being a Nazirite, he shall not eat\\
		anything that is made from the grapevine, from seeds to\\
		skin. {[image]}\\
	}
}
{
}

\Verse{5}
{
	\hP{כּל יְמֵי נֶדֶר נִזְרוֹ תַּעַר לֹא יַעֲבֹר עַל רֹאשׁוֹ עַד מְלֹאת הַיָּמִם אֲשֶׁר יַזִּיר לַיהֹוָה קָדֹשׁ יִהְיֶה גַּדֵּל פֶּרַע שְׂעַר רֹאשׁוֹ׃}
}
{
	\eP{
		All the days of his Nazirite vow, a razor shall not pass\\
		on his head; until the fulfillment of the days that he\\
		will make a separation for YHWH, he shall be holy, growing\\
		the hair of his head loose. {[image]}\\
	}
}
{
%	\footnote{
%		he shall be holy. This, too, is like priests (“They shall be holy to
%		their God,” Lev 21:6). And since it especially mentions the name of God
%		(“a separation for YHWH”), it is especially like the high priest (“I,
%		YHWH, make him holy,” Lev 21:15).
%	}
}

\Verse{6}
{
	\hP{כּל יְמֵי הַזִּירוֹ לַיהֹוָה עַל נֶפֶשׁ מֵת לֹא יָבֹא׃}
}
{
	\eP{
		All the days of his making a separation for YHWH he shall\\
		not come to a dead person. {[image]}\\
	}
}
{
%	\footnote{
%		come to a dead person. Meaning: touch a dead body. Interestingly, the
%		case of Samson is different. Like Samuel, he is a Nazirite from birth,
%		and he is forbidden to cut his hair—which becomes the key to his story.
%		But he is not forbidden to have contact with dead bodies—which is also
%		essential to the story, because he becomes a warrior and personally
%		makes thousands of dead bodies!
%	}
}

\Verse{7}
{
	\hP{לְאָבִיו וּלְאִמּוֹ לְאָחִיו וּלְאַחֹתוֹ לֹא יִטַּמָּא לָהֶם בְּמֹתָם כִּי נֵזֶר אֱלֹהָיו עַל רֹאשׁוֹ׃}
}
{
	\eP{
		For his father and for his mother, for his brother and for\\
		his sister: he shall not become impure for them when they\\
		die, because the crown of his God is on his head. {[image]}\\
	}
}
{
%	\footnote{
%		the crown of his God is on his head. As I noted in the case of the high
%		priest, the Hebrew term n[image]zer has the same range as the English
%		word crown, denoting either the top of the head or the headdress worn by
%		a king or priest (see my comment on Lev 21:12). Here it refers to the
%		hair on the Nazirite’s head, which may not be cut. As long as the man or
%		woman wears the uncut Nazirite hair, he or she also follows the high
%		priest’s practice of not touching a corpse.
%	}
}

\Verse{8}
{
	\hP{כֹּל יְמֵי נִזְרוֹ קָדֹשׁ הוּא לַיהֹוָה׃}
}
{
	\eP{
		All the days of his being a Nazirite he is holy to YHWH.\\
		{[image]}\\
	}
}
{
}

\Verse{9}
{
	\hP{וְכִי יָמוּת מֵת עָלָיו בְּפֶתַע פִּתְאֹם וְטִמֵּא רֹאשׁ נִזְרוֹ וְגִלַּח רֹאשׁוֹ בְּיוֹם טהֳרָתוֹ בַּיּוֹם הַשְּׁבִיעִי יְגַלְּחֶנּוּ׃}
}
{
	\eP{
		And if someone will die by chance suddenly by him and will\\
		make his head’s crown impure, then he shall shave his head\\
		on the day of his purification: he shall shave it on the\\
		seventh day. {[image]}\\
	}
}
{
}

\Verse{10}
{
	\hP{וּבַיּוֹם הַשְּׁמִינִי יָבִא שְׁתֵּי תֹרִים אוֹ שְׁנֵי בְּנֵי יוֹנָה אֶל הַכֹּהֵן אֶל פֶּתַח אֹהֶל מוֹעֵד׃}
}
{
	\eP{
		And on the eighth day he shall bring two turtledoves or\\
		two pigeons to the priest, to the entrance of the Tent of\\
		Meeting. {[image]}\\
	}
}
{
}

\Verse{11}
{
	\hP{וְעָשָׂה הַכֹּהֵן אֶחָד לְחַטָּאת וְאֶחָד לְעֹלָה וְכִפֶּר עָלָיו מֵאֲשֶׁר חָטָא עַל הַנָּפֶשׁ וְקִדַּשׁ אֶת רֹאשׁוֹ בַּיּוֹם הַהוּא׃}
}
{
	\eP{
		And the priest shall do one as a sin offering and one as a\\
		burnt offering and shall make atonement over him in that\\
		he has sinned over a person. And he shall make his head\\
		holy on that day. {[image]}\\
	}
}
{
%	\footnote{
%		a person. Meaning: a dead person.
%	}
}

\Verse{12}
{
	\hP{וְהִזִּיר לַיהֹוָה אֶת יְמֵי נִזְרוֹ וְהֵבִיא כֶּבֶשׂ בֶּן שְׁנָתוֹ לְאָשָׁם וְהַיָּמִים הָרִאשֹׁנִים יִפְּלוּ כִּי טָמֵא נִזְרוֹ׃}
}
{
	\eP{
		And he shall make a separation for YHWH for the days of\\
		his being a Nazirite, and he shall bring a lamb in its\\
		first year for a guilt offering, and the first days shall\\
		fall because his Nazirite state became impure. {[image]}\\
	}
}
{
%	\footnote{
%		the first days shall fall. The time that he spent as a Nazirite before
%		becoming impure by contact with a corpse does not count toward the time
%		that he vowed to spend as a Nazirite. It is nullified by the contact.
%	}
}

\Verse{13}
{
	\hP{וְזֹאת תּוֹרַת הַנָּזִיר בְּיוֹם מְלֹאת יְמֵי נִזְרוֹ יָבִיא אֹתוֹ אֶל פֶּתַח אֹהֶל מוֹעֵד׃}
}
{
	\eP{
		“And this is the instruction for the Nazirite: On the day\\
		of the fulfillment of the days of his being a Nazirite one\\
		shall bring him to the entrance of the Tent of Meeting.\\
		{[image]}\\
	}
}
{
}

\Verse{14}
{
	\hP{וְהִקְרִיב אֶת קרְבָּנוֹ לַיהֹוָה כֶּבֶשׂ בֶּן שְׁנָתוֹ תָמִים אֶחָד לְעֹלָה וְכַבְשָׂה אַחַת בַּת שְׁנָתָהּ תְּמִימָה לְחַטָּאת וְאַיִל אֶחָד תָּמִים לִשְׁלָמִים׃}
}
{
	\eP{
		And he shall make his offering to YHWH: one unblemished\\
		lamb in its first year for a burnt offering and one\\
		unblemished ewe-lamb in its first year for a sin offering\\
		and one unblemished ram for a peace offering {[image]}\\
	}
}
{
}

\Verse{15}
{
	\hP{וְסַל מַצּוֹת סֹלֶת חַלֹּת בְּלוּלֹת בַּשֶּׁמֶן וּרְקִיקֵי מַצּוֹת מְשֻׁחִים בַּשָּׁמֶן וּמִנְחָתָם וְנִסְכֵּיהֶם׃}
}
{
	\eP{
		and a basket of unleavened bread of fine flour, cakes\\
		mixed with oil and unleavened wafers with oil poured on\\
		them with their grain offering and their libations.\\
		{[image]}\\
	}
}
{
}

\Verse{16}
{
	\hP{וְהִקְרִיב הַכֹּהֵן לִפְנֵי יְהֹוָה וְעָשָׂה אֶת חַטָּאתוֹ וְאֶת עֹלָתוֹ׃}
}
{
	\eP{
		And the priest shall bring it forward in front of YHWH and\\
		shall do his sin offering and his burnt offering. {[image]}\\
	}
}
{
}

\Verse{17}
{
	\hP{וְאֶת הָאַיִל יַעֲשֶׂה זֶבַח שְׁלָמִים לַיהֹוָה עַל סַל הַמַּצּוֹת וְעָשָׂה הַכֹּהֵן אֶת מִנְחָתוֹ וְאֶת נִסְכּוֹ׃}
}
{
	\eP{
		And he shall make the ram a peace-offering sacrifice to\\
		YHWH with the basket of unleavened bread, and the priest\\
		shall do his grain offering and his libation. {[image]}\\
	}
}
{
}

\Verse{18}
{
	\hP{וְגִלַּח הַנָּזִיר פֶּתַח אֹהֶל מוֹעֵד אֶת רֹאשׁ נִזְרוֹ וְלָקַח אֶת שְׂעַר רֹאשׁ נִזְרוֹ וְנָתַן עַל הָאֵשׁ אֲשֶׁר תַּחַת זֶבַח הַשְּׁלָמִים׃}
}
{
	\eP{
		And the Nazirite shall shave his head’s crown at the\\
		entrance of the Tent of Meeting and shall take the hair of\\
		his head’s crown and put it on the fire that is under the\\
		peace-offering sacrifice. {[image]}\\
	}
}
{
}

\Verse{19}
{
	\hP{וְלָקַח הַכֹּהֵן אֶת הַזְּרֹעַ בְּשֵׁלָה מִן הָאַיִל וְחַלַּת מַצָּה אַחַת מִן הַסַּל וּרְקִיק מַצָּה אֶחָד וְנָתַן עַל כַּפֵּי הַנָּזִיר אַחַר הִתְגַּלְּחוֹ אֶת נִזְרוֹ׃}
}
{
	\eP{
		And the priest shall take the cooked shoulder from the ram\\
		and one unleavened cake from the basket and one unleavened\\
		wafer and put them on the Nazirite’s hands after his crown\\
		is shaved. {[image]}\\
	}
}
{
}

\Verse{20}
{
	\hP{וְהֵנִיף אוֹתָם הַכֹּהֵן תְּנוּפָה לִפְנֵי יְהֹוָה קֹדֶשׁ הוּא לַכֹּהֵן עַל חֲזֵה הַתְּנוּפָה וְעַל שׁוֹק הַתְּרוּמָה וְאַחַר יִשְׁתֶּה הַנָּזִיר יָיִן׃}
}
{
	\eP{
		And the priest shall elevate them, an elevation offering\\
		in front of YHWH. It is holy for the priest along with the\\
		breast of the elevation offering and the thigh of the\\
		donation. And after that the Nazirite may drink wine.\\
		{[image]}\\
	}
}
{
}

\Verse{21}
{
	\hP{זֹאת תּוֹרַת הַנָּזִיר אֲשֶׁר יִדֹּר קרְבָּנוֹ לַיהֹוָה עַל נִזְרוֹ מִלְּבַד אֲשֶׁר תַּשִּׂיג יָדוֹ כְּפִי נִדְרוֹ אֲשֶׁר יִדֹּר כֵּן יַעֲשֶׂה עַל תּוֹרַת נִזְרוֹ׃}
}
{
	\eP{
		This is the instruction for the Nazirite who will vow his\\
		offering to YHWH for his being a Nazirite outside of what\\
		his hand may attain. According to his vow that he will\\
		make, so he shall do along with the instruction of his\\
		being a Nazirite.” {[image]}\\
	}
}
{
}

\Verse{22}
{
	\hP{וַיְדַבֵּר יְהֹוָה אֶל מֹשֶׁה לֵּאמֹר׃}
}
{
	\eP{And YHWH spoke to Moses, saying, [image]}
}
{
}

\Verse{23}
{
	\hP{דַּבֵּר אֶל אַהֲרֹן וְאֶל בָּנָיו לֵאמֹר כֹּה תְבָרְכוּ אֶת בְּנֵי יִשְׂרָאֵל אָמוֹר לָהֶם׃}
}
{
	\eP{
		“Speak to Aaron and to his sons, saying, This is how you\\
		shall bless the children of Israel; say to them: {[image]}\\
	}
}
{
%	\footnote{
%		This is how you shall bless. A new form of expression that enters
%		biblical narrative in the book of Numbers is prayer. In earlier books
%		there have been depictions of individuals in prayer, but Numbers
%		includes a formal prayer, commanded to be said verbatim by the priests,
%		known as the priestly benediction (6:22–27). This blessing has been
%		pronounced upon the congregation by those of priestly descent presumably
%		from the biblical period until the present day (in some synagogues). It
%		has come to be of special interest recently. While staying in Jerusalem
%		in the summer of 1978, I looked out my window and watched a team of
%		archaeologists working below. They had uncovered Jewish tombs of the
%		Iron Age. I (and they) did not know it then, but inscriptions in thin
%		silver foil that they discovered were later painstakingly unrolled and
%		found to bear the words of the priestly benediction. The silver foil
%		inscriptions are the oldest known texts of a passage from the Bible.
%		(Their wording differs slightly from the biblical text and from each
%		other.)
%	}
}

\Verse{24}
{
	\hP{יְבָרֶכְךָ יְהֹוָה וְיִשְׁמְרֶךָ׃}
}
{
	\eP{May YHWH bless you and watch over you. [image]}
}
{
%	\footnote{
%		you. The language of the priestly benediction, like the Decalogue, is
%		formulated in the second-person singular. Although the priests are
%		instructed to bless the entire people with it, and although it
%		historically has been pronounced as a congregational blessing, it is
%		both a personal and a communal wish. In three brief lines the priests
%		wish each individual Israelite divine blessing, protection, favor,
%		grace, attention, and peace.
%	}
}

\Verse{25}
{
	\hP{יָאֵר יְהֹוָה פָּנָיו אֵלֶיךָ וִיחֻנֶּךָּ׃}
}
{
	\eP{
		May YHWH make His face shine to you and be gracious to\\
		you. {[image]}\\
	}
}
{
}

\Verse{26}
{
	\hP{יִשָּׂא יְהֹוָה פָּנָיו אֵלֶיךָ וְיָשֵׂם לְךָ שָׁלוֹם׃}
}
{
	\eP{May YHWH raise His face to you and give you peace. [image]}
}
{
%	\footnote{
%		and give you peace. The culminating wish of peace in the priestly
%		benediction is particularly poignant because Numbers is the first book
%		of the Bible so far to concentrate on war. Certainly peace can mean a
%		secure condition for an individual generally rather than being limited
%		to meaning an absence of war, but the issue of war is so much a part of
%		Numbers that the reference to peace in this significant position is
%		nonetheless quite striking. It appears in the beginning section of the
%		book as part of the preparation at Mount Sinai for the journey to come.
%		The people’s journey then comes to include substantial doses of reality
%		in the form of confrontations with hostile groups. Israel’s only
%		experience of war before this is in the brief account of an attack by
%		the Amalekites shortly after the exodus as the people first approach
%		Horeb (Exod 17:8–16). Indeed, the book of Exodus reports that the reason
%		the Israelites go to the Red Sea on the way out of Egypt instead of
%		taking the more obvious route, overland through Philistine territory, is
%		that YHWH is concerned that they are not yet ready to face war (13:17).
%		When the Egyptian army pursues them, Moses instructs them, “YHWH will
%		fight for you, and you’ll keep quiet!” (14:14); and again in the Song of
%		the Sea, it is the deity who is described as a warrior (15:3). Genesis
%		recounts only one battle, in which Abraham becomes involved because of
%		the capture of his nephew Lot (Genesis 14). Leviticus includes no
%		battles, and the word “war” does not occur in it. Deuteronomy, as I
%		shall discuss, treats the subject of war through remarkable legislation
%		but not through narrative other than to refer to the war episodes of
%		Numbers. It is Numbers alone among the Five Books of Moses that
%		especially develops this worst of human conditions through narrative of
%		specific military encounters. And it is Numbers that says, “May YHWH
%		give you peace”!
%	}
}

\Verse{27}
{
	\hP{וְשָׂמוּ אֶת שְׁמִי עַל בְּנֵי יִשְׂרָאֵל וַאֲנִי אֲבָרְכֵם׃}
}
{
	\eP{
		And they shall set my name on the children of Israel, and\\
		I shall bless them.” {[image]}\\
	}
}
{
}
